\section{Introduction}\label{sec:intro}

\begin{figure*}[t]
  \centering
  \includegraphics[width=0.90\textwidth]{20_snap_hyre_pipeline_art.png}
  \caption{Snap-HyRE (ours). Call 1 drafts privately and writes search text;
  only the search text is retrieved. Call 2 answers from the original question
  and retrieved evidence.}
  \label{fig:pipeline_art}
\end{figure*}

Legal RAG should work well for legal question answering: many answers turn on
statutes, regulations, holdings, or compact doctrinal rules, and retrieval can
put that text in front of a language model \citep{lewis2020rag}. The hard part
is often the first query. A bar-exam fact pattern may mention a driver,
crutches, and a banana peel, while the support passage states a black-letter
rule about proximate cause. A housing question may name a state, while an
unscoped national corpus can return a plausible statute from the wrong
jurisdiction. Final answer accuracy alone hides these failures.

Generated-query retrieval attacks this query mismatch. HyDE embeds a
hypothetical document rather than the user's original query \citep{gao2023hyde};
query expansion and rewriting use related generated search text
\citep{wang2023query2doc}. These methods are attractive in law because a model
can often infer the legal category behind a fact pattern before retrieval sees
the corpus. They can also fail by abstracting away anchors, choosing the wrong
legal frame, or retrieving plausible but irrelevant authority.

We introduce \hyre{}, a fixed two-call generated-query method. The first model
call writes two objects: a draft answer and a separate legal search passage.
The draft answer is private. It can guide the vocabulary of the search passage,
but it is not embedded, retrieved, or shown to the final answer call. The
second call answers the original question from the retrieved legal text. The
method uses model priors for query formation while keeping answer selection
tied to corpus evidence.

We report a two-benchmark study. BarExamQA tests the
weak-query regime, where fact-heavy questions need to be translated into
rule-shaped retrieval queries. HousingQA tests the corpus-interface regime,
where retrieval is not meaningful until every method searches only the
jurisdiction named in the question.

The main result is sharpest on BarExamQA. Across Llama 3.1 8B, Gemma 4 26B,
and Llama 3.3 70B, \hyre{} improves raw-question RAG by +2.4, +4.0, and +5.1
percentage points. It also raises BarExamQA Hit@5 from 1.4\% for raw-question
retrieval to 9.5--12.1\% for the three \hyre{} generators. HousingQA gives a
different lesson: adding the jurisdiction state filter raises raw Hit@5 from
2.8\% to 36.9\%, making state-filtered raw retrieval a strong baseline rather
than a method to beat by default.

This paper makes three contributions: it defines \hyre{} as a fixed two-call
legal RAG method; it compares raw-question RAG, HyDE, \hyre{}, and oracle
evidence controls on BarExamQA and HousingQA; and it analyzes when the method
helps by reporting both answer accuracy and gold evidence in the top five
retrieved passages. Structured query rewriting and additional diagnostics are
reported in the appendix.
